\chapter{The Virtual Fluid Network}
\label{ch:fluidnet}

What the virtual grid is for power, this is for fluid: it moves a fluid between places with no pipe
between them.

\begin{iefork}
Not in stock Immersive Engineering. It is deliberately built as a mirror of the power grid rather
than as a generalisation of it --- the two have separate save data, so a mistake in one cannot cost
you the other.
\end{iefork}

\section{Mains}

The fluid network is divided into \textbf{mains}, which are what segments are to the grid: named
groups of fittings with their own switch, limits and statistics.

A main carries \textbf{one fluid}. The first inlet to deliver decides which; after that the main
refuses anything else until it drains. That is not a limitation to work around --- it is what makes
a main something you can name ``diesel'' and trust.

\section{Fittings}

\begin{description}
\item[Fluid Inlet] Takes fluid from an adjacent tank or pipe and puts it into its main.
\item[Fluid Outlet] Draws from its main and fills whatever sits against it. One marked
      \textbf{critical} keeps its supply when the main runs short.
\item[Main Valve] Shuts a main by hand or on a redstone signal, and reports whether it is
      pressurised.
\item[Fluid Control Console] Where the network is managed. A separate structure from the power
      console, because the two networks fail apart.
\end{description}

\section{Failover and schedules}

As the grid: a main may name a backup, and may carry a schedule that can hold it closed but never
open one somebody shut.

\section{City mode}

Mains stop accounting for millibuckets and switch to presence --- a main is pressurised or it is
not.
